\documentclass[11pt]{article}
\usepackage{ochamath}
\subjectcolor{2F5DA8}
\setsheet{線形代数}{ベクトル空間と部分空間　演習}{https://university-math-crj.pages.dev/linear-algebra/vector-spaces/}
\begin{document}
\makesheethead{解答}

\problem[部分空間の判定]
\begin{ans}
$W_1=\{t(1,3):t\in\mathbb R\}$ は零を含み、和と実数倍で閉じているので部分空間。$W_2$ は $(0,0)$ を含まないため部分空間ではない。
\end{ans}
\point{含まれる点1つの検査だけでは、閉じていることの証明にならない。}

\problem[生成の範囲]
\begin{ans}
前者は $3\boldsymbol u+2\boldsymbol v=(3,2,6-2)^{\mathsf T}=(3,2,4)^{\mathsf T}$。一般の組合せは $(s,t,2s-t)^{\mathsf T}$ なので、生成する平面は $z=2x-y$。後者では最初の2成分から $s=t=0$ となり、第3成分も0になるため作れない。
\end{ans}
\point{成分ごとの方程式で生成に含まれるかを判定する。}
\newpage

\problem[独立性]
\begin{ans}
$2\boldsymbol a-\boldsymbol b=\boldsymbol0$ が非自明な関係なので従属。後者は $s(1,0)^{\mathsf T}+t(0,1)^{\mathsf T}=\boldsymbol0$ から $s=t=0$ のみなので独立。
\end{ans}
\point{非零ベクトルであることと、組が独立であることは別。}

\problem[出力制約]
\begin{ans}
$y_1=s,\ y_2=t,\ y_3=s+t$ なので $y_3=y_1+y_2$。$(3,2,5)^{\mathsf T}$ は $s=3,t=2$ で到達可能。
\end{ans}
\point{成分の数と、独立に変えられる数を区別する。}
\end{document}
